\section{碰撞偵測}
\label{sec:collision}

\subsection{碰撞形狀}

機器人的碰撞箱是一個直立圓柱，以腳底中心為基準；障礙物則依形狀使用可繞 $y$ 軸旋轉的直立方塊、直立圓柱或球。
這些形狀都不會傾斜，碰撞計算可以拆成鉛直與水平兩部分處理，比一般的三維多面體簡單許多。

圓柱的大小依姿勢而定：站立時半徑 0.45 m、高 1.75 m；潛行時身體前傾，高度降為 1.57 m，中心往前偏 0.11 m；坐下時雙腿前伸，高度只剩 1.06 m，中心往前偏 0.34 m。
這些數值依各姿勢下身體占據的範圍決定，偏移量隨身體朝向旋轉，轉換姿勢時也與關節角度一起內插，因此潛行時能鑽進較低的空間，坐下時腳也不會穿進障礙物；圖~\ref{fig:stances} 中的線框即為各姿勢的碰撞箱。

\subsection{偵測與推擠}

每個 tick 先依速度移動碰撞箱，再檢查與附近 $3 \times 3$ 個 chunk 中障礙物的重疊。
偵測函式回傳接觸法線 $\mathbf{n}$ 與穿透深度 $\delta$，三種障礙物的作法如下：
\begin{itemize}[nosep]
    \item \textbf{方塊}：把圓柱中心轉到方塊的局部座標，將中心夾在方塊的 $x$、$z$ 範圍內得到最近點，最近點到圓心的距離小於半徑即為碰撞。
    \item \textbf{圓柱}：兩圓心的水平距離小於半徑和即為碰撞。
    \item \textbf{球}：在圓柱上找離球心最近的點（高度夾在圓柱上下緣之間，水平方向夾在圓周內），距離小於球的半徑即為碰撞。
\end{itemize}
求出水平與鉛直的重疊量後，取較小的方向作為推出方向；例如從方塊上方落下時鉛直重疊較小，機器人會被往上推，站到方塊頂端。

碰撞箱沿 $-\mathbf{n}$ 推出 $\delta$，並移除速度中朝向障礙物的分量 $(\mathbf{v} \cdot \mathbf{n})\,\mathbf{n}$，保留沿著表面的分量，所以斜著撞牆時會順著牆面滑開，而不是整個停住。
推出後可能又撞到另一個障礙物，因此重複推擠，直到沒有重疊或最多 5 次為止。
表面剛好相切時也視為接觸（深度為 0），站在方塊頂端時才能穩定判定為著地。

\subsection{陡坡與支撐面}

地形以三角形組成，每個位置都可以查詢高度與法向量。
法向量與 $y$ 軸的夾角小於 $45^\circ$（$n_y > \frac{\sqrt{2}}{2}$）時視為可以行走；更陡的坡則當成牆壁：腳底低於地面時，沿坡面法向量的水平分量把機器人往外推，推出的距離為高度差乘上 $\cot$（坡度），因此機器人爬不上陡坡。

推擠完成後，在腳底稍微往下探，找出「支撐面」的高度：可行走的地形，以及頂面朝上的障礙物中最高的一個。
腳底與支撐面的距離為 0 時算著地；若上一個 tick 已經著地、沒有往上跳，且距離在 0.3 m 以內，也會直接貼回支撐面上。
沒有這項處理時，往下坡走的機器人每一步都會短暫離地，被當成在空中而無法控制方向。
